% Add "answers" option to print the answers
\documentclass[addpoints, answers]{exam}

\usepackage[utf8]{inputenc}
\usepackage{calc}
\usepackage{listings}
\usepackage{xcolor}
\usepackage{xspace}

\definecolor{codegreen}{rgb}{0,0.6,0}
\definecolor{codegray}{rgb}{0.5,0.5,0.5}
\definecolor{codepurple}{rgb}{0.58,0,0.82}

\lstdefinestyle{Cstyle}{
    commentstyle=\color{codegreen},
    keywordstyle=\color{magenta},
    numberstyle=\tiny\color{codegray},
    stringstyle=\color{codepurple},
    basicstyle=\ttfamily\small,
    breakatwhitespace=true,
    breaklines=true,
    captionpos=b,
    keepspaces=true,
    numbers=left,
    numbersep=5pt,
    showspaces=false,
    showstringspaces=false,
    showtabs=false,
    tabsize=4
}

\lstset{style=Cstyle}

% \makebox[\textwidth]{Name:\enspace\hrulefill}
%
% \vspace{5mm}
%\makebox[\textwidth]{Instructor’s name:\enspace\hrulefill}

\newcommand{\truefalse}
{		\begin{checkboxes}
			 \choice True
			 \choice False. Explanation:
		\end{checkboxes}
		\vspace{10pt}
}

\newcommand{\answerspace}[1]
{
	\ifprintanswers
	\else
		\setlength{\mylength}{10pt}%
		%\noindent\rule{\mylength}{20pt}
		%\vspace{\fpeval{1.2*(#1)}}
		\vspace{#1\mylength}
	\fi
}
\newlength{\mylength}

\newcommand\pf[1]{\textcolor{red}{PF: #1}}
\newcommand\am[1]{\textcolor{blue}{AM: #1}}

\newcommand{\answer}[1]{\begin{solution}{#1}\end{solution}}
\newcommand{\ra}{$\rightarrow$\xspace}


\usepackage{color}
\usepackage{lastpage}
\usepackage{verbatim}
\usepackage{fancyvrb}
\usepackage{graphicx}
\usepackage{float}
\usepackage{minted}
\usepackage{multicol}
\usepackage{amsmath}
\usepackage[normalem]{ulem}

\usepackage[svgnames]{xcolor}
\usepackage[colorlinks=true, allcolors=MidnightBlue]{hyperref}

% For those sweet, sweet acmart fonts.
\usepackage[T1]{fontenc}
\usepackage[tt=false, type1=true]{libertine}
\usepackage[varqu]{zi4}                  % Inconsolata for \ttfamily

% Eliminate extraneous spacing on itemize, description, etc.
\usepackage{enumitem}
\setlist{nosep}

\newcommand{\Code}[1]{\mintinline{python}{#1}}
\newcommand{\nonAnswerNewPage}{\ifprintanswers \vfill \else \newpage \fi}

\renewcommand{\solutiontitle}{\noindent\textbf{Solution sketch: }}
\definecolor{SolutionColor}{rgb}{0.8,0.9,1}

\cfoot{Page \thepage\ of \pageref{LastPage}}
\rhead{Initials: \rule{20mm}{0.4pt}}

%%%%%%%%%%%%%%%%%%%%%%%%%%%%%%%%%%%%%%%%
\begin{document}
\fvset{numbers=left,numbersep=3pt}
\thispagestyle{empty}

\begin{center}
  {\LARGE{\textbf{CS 408: Homework 4}}}\\
  \vspace*{5pt}
  {\Large Deadline: Friday, October 16th \textbf{at 6PM}}\\
  \vspace*{10pt}

  \fbox{
    \fbox{\parbox{6in}{\centering
        \begin{itemize}
        \item The homework is open notes, open book.
        \item Please upload an annotated version of this pdf, ensuring
          that each answer fits in the box following each question.
        \item All solutions must be \textbf{handwritten}; typed work
          will not receive any credit. You may annotate the pdf
          directly, no need to print out and rescan.
        \item No AI tools allowed.
        \item If a question asks for code, you do not need to
          reproduce Hypothesis's syntax exactly, but make sure that we
          can clearly identify what each strategy produces and what
          each property is testing.
        \item No collaboration with your peers of any kind allowed on
          this homework.
        \item Homeworks are due at \textbf{6PM}. Anything submitted
          after \textbf{6PM} will not be accepted.
        \item Provide complete and clear but concise answers. Clearly
          explain your answers where required.
	\end{itemize}
      }}}
\end{center}

\begin{center}
	\makebox[0.6\textwidth]{
		\combinedgradetable[h][questions]
	}
\end{center}
%\begin{center}\combinedgradetable[h][questions]\end{center}

\begin{center}
\makebox[\textwidth]{\textbf{Name (please print)}: \rule{100mm}{0.8pt} }
\end{center}

\vspace{10mm}

\begin{questions}

%% =====================================================================
%% Q1 -- test oracles
%% =====================================================================

\question
  A team is building a \Code{spell-checker} function that
  takes as input a word, and returns a list of suggested corrections,
  drawn from the \href{https://www.oed.com}{Oxford English Dictionary}
  (OED), ordered from ``best'' to ``worst''. The team wants to test
  \Code{spell-checker} automatically.

  Explain in 1--2 sentences why each of the following three test
  oracles is not a good fit for this testing scenario.
  %
  \begin{parts}
\part[2] An \textbf{implicit} test oracle that checks a) whether the
  program crashed and b) whether the words in the output list are
  included in the OED.
  \begin{solutionorbox}[2in]
    Returning \emph{a} list of suggestions without crashing says
    nothing about the quality of those suggestions: even a bad answer
    is likely to be a well-formed list of words.
  \end{solutionorbox}
  \nonAnswerNewPage
\part[2] Asking two junior engineers to act as \textbf{human} oracles.
  \begin{solutionorbox}[2in]
    A person can only judge a handful of suggestions, and they are
    unlikely to scan the entirety of the OED when checking the
    word. Moreover, people will rank the same suggestions
    differently, and we have no way of adjudicating disagreements.
  \end{solutionorbox}
\part[2] Implementing and using a \textbf{reference} implementation.
  \begin{solutionorbox}[2in]
    This requires us to implement a second spell-checker function that
    we trust to be correct. Unlike sorting, however, there is no
    slow-but-obviously-correct algorithm to use as ground truth,
    especially since the quality of solutions is a judgement call.
\end{solutionorbox}
\end{parts}

\question[6] Give \textbf{three properties} you could check about
  \Code{spell-checker}'s output that do not require knowing the
  correct answer. At least one should relate the outputs of two
  \emph{different} inputs. %
  \begin{solutionorbox}[3.2in]
    There are many correct answers here:
    \begin{itemize}
    \item every suggestion is a word that actually appears in the OED;
    \item no suggestion appears twice;
    \item if the word already appears as-is in the dictionary, it is
      suggested first, as it (likely) spelled correctly;
    \item all the output words are within some edit distance of the
      input word.
    \end{itemize}
    Relations between two answers, i.e., metamorphic relations:
    \begin{itemize}
    \item changing the case of the misspelled word does not change the
      suggestions;
    \item adding an unrelated word to the OED does not change the
      relative order of the existing suggestions;
    \item running the same query twice returns the same list, i.e.,
      the function is deterministic.
    \end{itemize}
  \end{solutionorbox}

  \nonAnswerNewPage
\question[1] None of your answers to the second question establishes
  that a suggestion is \emph{fully correct}. Why is this an acceptable
  tradeoff here? %
  \begin{solutionorbox}[2in]
    As the answer to 1c explains, there is no obvious definition of
    correct here. Using weaker properties that any reasonable
    implementation should satisfy will (hopefully) catch some faults,
    and using relations in particular allows us to ensure that
    \Code{spell-checker} is internally consistent. In short: an
    imperfect oracle that we can actually use is better than a perfect
    one that we either do not have time or do not know how to
    implement.
  \end{solutionorbox}

%% =====================================================================
%% Q2 -- properties in Hypothesis
%% =====================================================================

\question Consider a function \Code{merge(xs, ys)} that merges two
  \textbf{sorted} lists of integers into a single sorted list, keeping
  duplicates.

\begin{parts}

\part[6] Write \textbf{three} Hypothesis properties for \Code{merge}.
  Each should be a separate test that is annotated with a
  \Code{@given} which explicitly states the input generation
  strategy. If necessary, you can define your own custom strategy,
  e.g.,
\begin{minted}[linenos=false]{python}
 nonzero_gen = st.integers().filter(lambda x: x != 0)
\end{minted}
%
  \begin{solutionorbox}[4.5in]
    Many correct answers. Note that since the inputs to \Code{merge}
    are assumed to be sorted, the properties need to either use an
    input generation strategy that always produces sorted lists (see
    line 1 for an example), or use \Code{assume} to filter them
    somehow.
\begin{minted}{python}
sorted_lists = st.lists(st.integers()).map(sorted)

@given(sorted_lists, sorted_lists)
def test_output_is_sorted(xs, ys):
    assert is_sorted(merge(xs, ys))

@given(sorted_lists, sorted_lists)
def test_length_is_preserved(xs, ys):
    assert len(merge(xs, ys)) == len(xs) + len(ys)

@given(sorted_lists, sorted_lists)
def test_same_elements(xs, ys):
    assert Counter(merge(xs, ys)) == Counter(xs) + Counter(ys)

@given(sorted_lists, sorted_lists)
def test_merge_is_commutative(xs, ys):
    assert merge(xs, ys) == merge(ys, xs)

@given(sorted_lists)
def test_merging_with_empty(xs):
    assert merge(xs, []) == xs
\end{minted}
  \end{solutionorbox}

\nonAnswerNewPage

\part[3] Which of your properties are \textbf{metamorphic relations},
  and which are not? Explain what distinguishes the two categories. %
  \begin{solutionorbox}[3in]
    A metamorphic relation constrains how the outputs of \emph{two
    related inputs} relate to one another; the others constrain a
    single output on its own.

    \textbf{Metamorphic:} \Code{test_merge_is_commutative} relates
    \Code{merge(xs, ys)} to \Code{merge(ys, xs)}. \\
    \Code{test_merging_with_empty} is a borderline case, since it
    relates the output to the input rather than to another output. It
    is better described as an identity.

    \textbf{Not metamorphic:} sortedness, length, and the multiset of
    elements are all properties of one output produce by a single pair
    of inputs.

    Note that none of the properties computes the correct answer. A
    reference implementation would, but then we would be testing
    \Code{merge} against something we already have.
  \end{solutionorbox}

\end{parts}

\nonAnswerNewPage

%% =====================================================================
%% Q3 -- a custom generator
%% =====================================================================

\question In Python, suppose we represent a color as one of the
  strings \Code{"red"}, \Code{"green"}, or \Code{"blue"}, and an
  \emph{optional} color as either such a string or \Code{None}.

\begin{parts}

\part[8] Write a Hypothesis generator for optional colors that
  produces \Code{None} (roughly) ten percent of the time, and produces
  \Code{"red"} (roughly) three times as often as either of the other
  two colors. %
  \begin{solutionorbox}[4in]
    Unlike some other PBT frameworks, e.g., QuickCheck, Hypothesis has
    no built-in weighted choice, so we have to be explicit about how
    values are weighted. One way to do this is by branching on a
    randomly \Code{draw}n number; another is to sample from a pool
    with repeated entries.
\begin{minted}{python}
@st.composite
def optional_color(draw):
    if draw(st.integers(min_value=1, max_value=10)) == 1:
        return None
    return draw(st.sampled_from(
        ["red", "red", "red", "green", "blue"]))
\end{minted}
    The first draw gives \Code{None} with probability
    $\frac{1}{10}$. The remaining $\frac{9}{10}$ is divided among a
    pool of options that contains three copies of \Code{"red"}, and
    thus will be chosen with probability
    $\frac{9}{10}\cdot\frac{3}{5} = \frac{27}{50}$, while the other
    two colors will be picked with probability
    $\frac{9}{10}\cdot\frac{1}{5} = \frac{9}{50}$, i.e., three times
    less than \Code{"red"}, as required.
  \end{solutionorbox}

\part[4] How does the following generator differ from the one you
  wrote?
  \begin{minted}[linenos=false]{python}
  st.one_of(st.none(), st.sampled_from(["red","green","blue"]))
\end{minted}
  \begin{solutionorbox}[2.5in]
    \Code{st.one_of} returns one of each of its arguments with equal
    probability, so it would thus produce \Code{None} roughly half the
    time. \Code{st.sampled_from} similarly returns one of the three
    colors with equal likelihood, producing \Code{red} about the same
    number of times as the other two colors, which is not what the
    question asks for.

    In general, a generator encodes a probability distribution on
    inputs to the system under test, and we may have to explicitly
    construct that distribution to ensure we're effectively sampling
    from that distribution.
  \end{solutionorbox}

% \part[3] Suppose you ran your generator 1000 times and counted how
%   often each value appeared. Would you expect the counts to match the
%   proportions above? Explain. %
%   \begin{solutionorbox}[2.5in]
%     Not if you count them through Hypothesis. There are only four
%     distinct values, and Hypothesis does not keep generating examples
%     it has already tried --- so a 1000-example run will produce only a
%     handful of actual draws, and the observed counts say nothing about
%     the distribution.

%     To check the distribution you would call the underlying random
%     choices directly, outside Hypothesis.
%   \end{solutionorbox}

\end{parts}

\nonAnswerNewPage

%% =====================================================================
%% Q4 -- mutation-based fuzzing
%% =====================================================================

\question Consider the \Code{cgi_decode} function below, which decodes
  a CGI-encoded string: \Code{"+"} becomes a space, \Code{"\%XX"}
  becomes the character with hex code \Code{XX}, and anything else is
  copied through.

\begin{minted}{python}
def cgi_decode(s):
  hex_values = {'0': 0, '1': 1, ..., 'e': 14, 'f': 15, 'A': 10, 'B': 11, ..., 'F': 15}
  t = ""
  i = 0
  while i < len(s):
    c = s[i]
    if c == '+':
      t += ' '                                 # line A
    elif c == '%':
      digit_high, digit_low = s[i+1], s[i+2]
      i += 2
      if digit_high in hex_values and digit_low in hex_values:
        t += chr(hex_values[digit_high] * 16 + hex_values[digit_low])   # line B
      else:
        raise ValueError("Invalid encoding")   # line C
    else:
      t += c                                   # line D
    i += 1
  return t
\end{minted}

  A mutation-based fuzzer starts from the seed population
  \{\ \Code{"keep"},\ \Code{"snore"},\ \Code{"unwise"},\ \Code{"quip"}\ \}
  and repeatedly picks a string from its pool, applies \textbf{one}
  mutation, and adds the result back to the pool. The three mutations
  are equally likely:
  \begin{description}
  \item[delete] remove the character at a uniformly chosen position;
  \item[insert] insert a uniformly chosen printable character
    (\Code{chr(32)}--\Code{chr(126)}, so 95 possibilities) at a
    uniformly chosen position;
  \item[flip] at a uniformly chosen position, replace the character
    \Code{c} with \Code{chr(ord(c) ^ (1 << b))} for a uniformly chosen
    bit $b \in \{0,\ldots,6\}$.
  \end{description}

  \noindent To illustrate how the last mutation works, observe that
  flipping the second bit of \Code{'u'} in \Code{"unwise"} gives
  \Code{"wnwise"}: \Code{'u'} is \Code{1110101} in binary, and
  flipping the second bit gives \Code{1110111}, which is
  \Code{'w'}. Importantly, this means that not every character is
  reachable from every other in a single flip. If you're aiming for a
  particular test input, it may be helpful to write the source and
  target character in binary.

\begin{parts}

\part[4] Which of lines A--D does each seed cover? Is this set of
  seeds a good starting point when fuzzing \Code{cgi_decode}? %
  \begin{solutionorbox}[2.25in]
    Every seed covers \textbf{line D only}, since they are all plain
    text, and do not contain either \Code{'+'} or \Code{'\%'}. From a
    coverage point of view, each of these starting seeds is
    equivalent.
  \end{solutionorbox}

\nonAnswerNewPage

\part[8] Give a set of inputs that achieves 100\% line coverage of
  \Code{cgi_decode}, and for each one give a \emph{sequence} of
  mutations that produces it from one of the seeds. Be specific about
  the position and the bit. %
  \begin{solutionorbox}[4.5in]
    Here is one possible set of inputs, each starting from the
    \Code{"keep"} seed:
    \begin{description}
    \item[line A] \Code{"+eep"} --- flip position 0:
      \Code{'k'} is \Code{0x6B} and \Code{'+'} is \Code{0x2B}, which
      differ in bit 6. \textbf{(1 mutation)}
    \item[line C] \Code{"k\%ep"} --- flip position 1:
      \Code{'e'} is \Code{0x65} and \Code{'\%'} is \Code{0x25}, again
      bit 6. The escape \Code{\%ep} is invalid, since \Code{'p'} is not
      a hex digit. \textbf{(1 mutation)}
    \item[line B] \Code{"k\%e0"} --- from \Code{"k\%ep"}, flip
      position 3: \Code{'p'} is \Code{0x70} and \Code{'0'} is
      \Code{0x30}, bit 6 again. Now \Code{\%e0} is a valid escape.
      \textbf{(2 mutations)}
    \end{description}
    Line D is already covered by every seed. Note that line B needs
    \emph{two} mutations: the \Code{'\%'} has to appear first, and
    there is no single bit flip that turns \Code{"keep"} into a string
    with a valid escape.
  \end{solutionorbox}

\nonAnswerNewPage

\part[6] For each of the inputs you included in your test set, what is
  the probability that a fuzzer would produce that input in a minimal
  number of runs? \sout{Assume that a mutant replaces its parent in
    each run, so that the seed pool always contains four elements.}
  Explain your reasoning for each. \sout{You can treat each input
    separately, starting from the original pool.}  %
  \begin{solutionorbox}[4in]
    \begin{itemize}
    \item \textbf{\Code{"+eep"} (line A)} --- one mutation from
      \Code{"keep"}: flip bit 6 at position 0. The probability is:
    \[
      \underbrace{\tfrac14}_{\text{pick the parent}} \cdot
      \underbrace{\tfrac13}_{\text{choose \emph{flip}}} \cdot
      \underbrace{\tfrac14}_{\text{choose the position}} \cdot
      \underbrace{\tfrac17}_{\text{choose the bit}} \;=\;
      \tfrac{1}{336}
    \]

  \item \textbf{\Code{"k\%ep"} (line C)} --- one mutation from
    \Code{"keep"}: flip bit 6 at position 1. The pool has one extra
    element now, so the probability is slightly less than the
    previous:
    $\tfrac15 \cdot \tfrac13 \cdot \tfrac14 \cdot \tfrac17 \;=\;
    \tfrac{1}{420}$

  \item \textbf{\Code{"k\%e0"} (line B)} --- one mutation using
    \Code{"k\%ep"}, which is now included in our pool of six seed
    strings:
    \[
      \tfrac16 \cdot \tfrac13 \cdot \tfrac14 \cdot \tfrac17
      \;=\; \tfrac{1}{504}
    \]
    flipping bit 6 at position 3 of \Code{"k\%ep"}.
  \end{itemize}

  \textbf{An aside:} Taking these all together, the likelihood of
  generating these inputs in exactly this sequence is:
    \[
      \tfrac{1}{336} \cdot \tfrac{1}{420} \cdot \tfrac{1}{504} \;=\;
      \tfrac{1}{71124480} \;\approx\; 1.4 \times 10^{-8}
    \]

    This assumes the fuzzer picks uniformly among the seeds when
    choosing the next one to mutate, and only considers this exact
    sequence of choices. If we only require that the chain happens
    \textit{eventually}, rather than immediately, the probability of
    generating these inputs is much higher.
  \end{solutionorbox}

\part[1] For a real-world fuzzer, why is it not a problem if the
  probability of generating a particular test input is small? %
  \begin{solutionorbox}[2.5in]
    On the one hand, real world fuzzers run for millions of
    iterations, so even a one-in-a-million event is likely to occur at
    least once.

    More importantly, a \emph{coverage-guided} fuzzer does not need
    this exact chain. When \Code{"k\%ep"} turns up it covers a new
    line, and the fuzzer keeps it and (potentially preferentially)
    mutates it further. Thus, a grey-box fuzzer does not blindly
    search the input space, starting from scratch at each iteration:
    useful mutations are kept and reused as the fuzzer keeps going.
  \end{solutionorbox}

\end{parts}

\end{questions}

\end{document}
